\chapter{Evidence\\สร้างความน่าเชื่อถือด้วยหลักฐาน}

ทุก Portfolio มีข้อกล่าวอ้าง ไม่ว่าจะเป็น \enquote{ระบบแม่นยำ} \enquote{ผู้ใช้พึงพอใจ} หรือ \enquote{งานสร้างผลกระทบ} Evidence คือสิ่งที่ทำให้ข้อกล่าวอ้างเหล่านั้นมีน้ำหนัก บทบาทของ Mentor ไม่ใช่ผลักให้ Mentee สะสมตัวเลขมากที่สุด แต่ช่วยเลือกหลักฐานที่เหมาะกับคำถามและการตัดสินใจของผู้เกี่ยวข้อง

\learningobjectives{
  \item แยกข้อกล่าวอ้างออกจากหลักฐานที่รองรับ
  \item สร้าง evidence ladder ที่เหมาะกับระยะของงาน
  \item ระบุข้อจำกัดและรักษาความซื่อตรงของข้อมูล
}

\section{เริ่มจากข้อกล่าวอ้าง}
ก่อนถามว่า \enquote{มีข้อมูลอะไร} ให้ถามว่า \enquote{ต้องการให้ใครเชื่ออะไรเพื่อทำการตัดสินใจใด} หากข้อกล่าวอ้างคือระบบตรวจโรคแม่นยำ หลักฐานควรเป็นผลทดสอบเทียบ baseline; หากข้อกล่าวอ้างคือเกษตรกรนำไปใช้ได้ หลักฐานจากห้องปฏิบัติการเพียงอย่างเดียวยังไม่พอ

\section{Evidence Ladder}
\begin{table}[ht]
\centering
\begin{tabularx}{\textwidth}{P{.23\textwidth}YP{.28\textwidth}}
\toprule
\textbf{ระดับ} & \textbf{ตัวอย่าง} & \textbf{ตอบคำถาม}\\
\midrule
แนวคิด & เหตุผลทางทฤษฎี ความเห็นผู้เชี่ยวชาญ & เป็นไปได้หรือไม่\\
ต้นแบบ & ผลทดสอบ lab, benchmark & ทำงานภายใต้เงื่อนไขควบคุมหรือไม่\\
การใช้จริง & pilot, observation, usage log & ผู้ใช้ใช้ได้จริงหรือไม่\\
ผลลัพธ์ & ก่อน--หลัง, ตัวชี้วัดต้นทุน/เวลา/คุณภาพ & เกิดการเปลี่ยนแปลงหรือไม่\\
การยืนยันภายนอก & audit, การรับรอง, การทำซ้ำโดยผู้อื่น & เชื่อถือและขยายผลได้หรือไม่\\
\bottomrule
\end{tabularx}
\end{table}

งานระยะแรกไม่จำเป็นต้องมีหลักฐานระดับสูงสุด แต่ควรซื่อตรงว่าหลักฐานปัจจุบันตอบคำถามใดและยังตอบอะไรไม่ได้ การระบุข้อจำกัดเพิ่มความน่าเชื่อถือมากกว่าการใช้ภาษาที่เกินข้อมูล

\begin{examplebox}
\textbf{ข้อกล่าวอ้างเดิม:} \enquote{ระบบช่วยลดเวลาตรวจอย่างมาก}\\
\textbf{Evidence ที่ชัดขึ้น:} \enquote{ในการทดลองกับเจ้าหน้าที่ 6 คน ระบบลดเวลามัธยฐานต่อชุดตัวอย่างจาก 18 นาทีเหลือ 11 นาที ภายใต้เงื่อนไข pilot สองสัปดาห์ ผลยังไม่รวมช่วงอบรมผู้ใช้ใหม่}\\
ข้อความหลังสื่อสารทั้งผล ขนาดตัวอย่าง บริบท และข้อจำกัด
\end{examplebox}

\section{Evidence Map}
ให้เชื่อมทุกข้อกล่าวอ้างสำคัญกับหลักฐานที่มี เจ้าของข้อมูล สถานที่จัดเก็บ คุณภาพ และหลักฐานที่ยังต้องสร้าง Evidence Map ช่วยให้การเก็บข้อมูลเป็นส่วนหนึ่งของ Process แทนที่จะย้อนกลับไปหาเมื่อถึงกำหนดส่งข้อเสนอ

\diagnostic{หลักฐานชิ้นใดจะลดความไม่แน่นอนของผู้ตัดสินใจได้มากที่สุด และเราสร้างมันอย่างมีจริยธรรมได้อย่างไร}

\begin{mentornote}
Evidence ที่มีตัวเลขไม่จำเป็นต้องแข็งแรงกว่าเรื่องเล่าเสมอไป ข้อมูลเชิงคุณภาพอธิบายกลไกและประสบการณ์ได้ดี เพียงต้องเก็บอย่างมีระบบและไม่เลือกเฉพาะเสียงที่สนับสนุนข้อสรุป
\end{mentornote}

\begin{keytakeaway}
Evidence ที่ดีสัมพันธ์โดยตรงกับข้อกล่าวอ้าง โปร่งใสเรื่องวิธีและข้อจำกัด และเพียงพอสำหรับการตัดสินใจในระยะนั้นของงาน
\end{keytakeaway}

\begin{reflection}
เขียนข้อกล่าวอ้างสำคัญสามข้อจาก Portfolio ระบุ Evidence ปัจจุบัน ระดับความแข็งแรง และหลักฐานชิ้นถัดไปที่ควรสร้าง พร้อมผู้รับผิดชอบและกำหนดเวลา
\blankline\blankline
\end{reflection}

\chaptersummary{Evidence เปลี่ยนความเชื่อให้เป็นข้อเสนอที่ตรวจสอบได้ การเริ่มจากข้อกล่าวอ้างและการตัดสินใจของผู้ใช้หลักฐาน ทำให้ Mentee เก็บข้อมูลอย่างมีจุดหมายและสื่อสารอย่างซื่อตรง}
